\section*{第 \ref{ch:b3:inorganic-materials} 章——无机材料与纳米材料}

\begin{solution}{exo:b3:inorganic-materials:1}
$x^2 \propto t$：厚度为三倍时，所需时间为九倍，即 \qty{90}{h}。
\end{solution}

\begin{solution}{exo:b3:inorganic-materials:2}
水解：$\ce{Si(OC2H5)4 + 4H2O -> Si(OH)4 + 4C2H5OH}$。\\
缩合：$\ce{2Si(OH)4 -> Si2O(OH)6 + H2O}$。\\
总反应：$\ce{Si(OC2H5)4 + 2H2O -> SiO2 + 4C2H5OH}$。
\end{solution}

\begin{solution}{exo:b3:inorganic-materials:3}
网络形成体为 $\ce{SiO2}$、$\ce{B2O3}$；网络修饰体为 $\ce{Na2O}$、$\ce{CaO}$。$\ce{Na2O}$ 的每个氧离子都会打断一个 Si--O--Si
桥键断裂为两个 Si--O$^-$ 端基，各与一个 $\ce{Na+}$ 配对：网络被切断，熔体
更易流动，熔化温度也更低。
\end{solution}

\begin{solution}{exo:b3:inorganic-materials:4}
十二个五边形（与所有富勒烯相同）。$V = 70$，$E = 3V/2 = 105$ 条键，$F = 2 - V + E = 37$
个面，故六边形有 $37 - 12 = 25$ 个。
\end{solution}

\begin{solution}{exo:b3:inorganic-materials:5}
$\sqrt2\,(r_B + r_O) = 1.4142 \times \qty{200.5}{pm} = \qty{283.5}{pm}$。$\ce{SrTiO3}$：$284/283.5 = 1.00$，为理想立方钙钛矿结构。
$\ce{BaTiO3}$：$301/283.5 = 1.06$，钛原子对其八面体而言过小，偏离中心位置：铁电体。
$\ce{CaTiO3}$：$274/283.5 = 0.97$，钙原子相对其空隙过小：八面体发生倾斜，
对称性降低。
\end{solution}

\begin{solution}{exo:b3:inorganic-materials:6}
Si/Al $= 106/86 = 1.23$。$M = 86 \times 23.0 + 86 \times 27.0 + 106 \times 28.1 + 384 \times 16.0 = \qty{13422.6}{g/mol}$；每个
化学式含 43 个 $\ce{Ca^2+}$，$43/\qty{13422.6}{g/mol} = \qty{3.20}{mmol/g}$。
\end{solution}

\begin{solution}{exo:b3:inorganic-materials:7}
最近邻距离 $d = a/\sqrt2 = \qty{288.4}{pm}$。含 $n$ 层壳层的团簇直径为 $(2n + 1)d$：
$2n + 1 \approx 5/0.2884 = 17.3$，故 $n = 8$（\qty{4.9}{nm}）。$N(8) = 2057$ 个原子，其中位于表面的有 $10 \times 64 + 2 = 642$ 个，占
31\,\%.
\end{solution}

\begin{solution}{exo:b3:inorganic-materials:8}
以 $h^2/(8\mu R^2)$ 计算，其中 $\mu = \qty{9.109e-32}{kg}$：在 \qty{2.0}{nm} 处为 \qty{0.94}{eV}，在 \qty{3.0}{nm} 处为 \qty{0.42}{eV}。发光能量为
$1.74 + 0.94 = \qty{2.68}{eV}$，即 \qty{463}{nm}（蓝光），$1.74 + 0.42 = \qty{2.16}{eV}$，即 \qty{574}{nm}（黄光）：较小的量子点
发出的光更偏蓝。
\end{solution}

\begin{solution}{exo:b3:inorganic-materials:9}
$(10, 10)$：扶手椅型，$n - m = 0$，金属性。$(10, 0)$：锯齿型，10 不是 3 的倍数，半导体性。
$(9, 0)$：锯齿型，金属性。$(12, 8)$：既非扶手椅型也非锯齿型（手性），$n - m = 4$，半导体性。
\end{solution}

\begin{solution}{exo:b3:inorganic-materials:10}
$M = 4 \times 65.4 + 13 \times 16.0 + 24 \times 12.0 + 12 \times 1.0 = \qty{769.6}{g/mol}$; $a^3 = (\qty{25.82e-8}{cm})^3 = \qty{1.721e-20}{cm^3}$;
$\rho = 8 \times 769.6/(\num{6.022e23} \times \num{1.721e-20}) = \qty{0.594}{g.cm^{-3}}$。足球场面积为 \qty{7140}{m^2}，一克沸石
的表面积即 $3000/7140 = 0.42$ 个足球场，一茶匙（数克）则超过一整个足球场。
\end{solution}

\begin{solution}{exo:b3:inorganic-materials:11}
$\dfrac{10n^2 + 2}{(10n^3 + 15n^2 + 11n + 3)/3} = \dfrac{30n^2 + 6}{10n^3 + 15n^2 + 11n + 3} \to \dfrac{30n^2}{10n^3} = \dfrac3n$。当 $n = 8$ 时，精确值
为 $642/2057 = 0.31$，对比 $3/8 = 0.375$；该近似高估了小团簇，
这些团簇的内层壳层不可忽略。
\end{solution}

\begin{solution}{exo:b3:inorganic-materials:12}
$2E = 3V$，$2E = 5p + 6h + 7s$，$F = p + h + s$。由欧拉公式 $V - E + F = 2$，代入 $V = 2E/3$ 得 $F - E/3 = 2$，于是
$p + h + s - (5p + 6h + 7s)/6 = 2$，即 $(p - s)/6 = 2$，$p - s = 12$。每个七边形需要
一个额外的五边形。
\end{solution}

\begin{solution}{pb:b3:inorganic-materials:1}
\textbf{1.} $\mathrm{Na_{12}Al_{12}Si_{12}O_{48}}$: $12 \times 23.0 + 12 \times 27.0 + 12 \times 28.1 + 48 \times 16.0 = \qty{1705.2}{g/mol}$.

\textbf{2.} $1705.2 + 27 \times 18.0 = \qty{2191.2}{g/mol}$.

\textbf{3.} Si/Al $= 12/12 = 1$.

\textbf{4.} 不存在 Al--O--Al 连接；Si/Al $= 1$ 时，每个 Si 有四个 Al 近邻，每个
Al 有四个 Si 近邻：严格交替排列。

\textbf{5.} $-12$，每个铝原子对应一个。

\textbf{6.} $486.0/2191.2 = 22.2$\,\%.

\textbf{7.} $\ce{2Na+}(\text{沸石}) + \ce{Ca^2+}(\text{aq}) \rightleftharpoons \ce{Ca^2+}(\text{沸石}) + \ce{2Na+}(\text{aq})$.

\textbf{8.} 六个（共十二个电荷）。

\textbf{9.} $6/\qty{1705.2}{g/mol} = \qty{3.52}{mmol/g}$.

\textbf{10.} $6/\qty{2191.2}{g/mol} = \qty{2.74}{mmol/g}$.

\textbf{11.} 每克相当于 $\qty{3.52}{mmol/g} \times \qty{100.1}{g/mol} = \qty{352}{mg}$ 的 $\ce{CaCO3}$。

\textbf{12.} $\qty{2.0}{mmol/L} \times \qty{15}{L} = \qty{30}{mmol}$.

\textbf{13.} $\qty{30}{mmol}/\qty{2.74}{mmol/g} = \qty{11}{g}$.

\textbf{14.} $\qty{10.9}{g}/0.60 = \qty{18}{g}$.

\textbf{15.} $\qty{60}{mmol}$ 的 $\ce{Na+}$，$\qty{0.060}{mol} \times \qty{23.0}{g/mol} = \qty{1.4}{g}$。

\textbf{16.} 废水中的磷酸盐使河流和湖泊中的藻类大量繁殖，
形成的水华及其腐败过程消耗溶解氧（富营养化）。

\textbf{17.} 较小的 $\ce{Mg^2+}$ 与水分子结合更紧密；它必须脱去水分子
才能通过窗口到达位点，而在洗涤温度下这一过程较慢。

\textbf{18.} 裸离子直径约为 \qty{0.20}{nm}，为窗口尺寸的一半。

\textbf{19.} 水合离子大于窗口，它必须脱去部分
水合层才能通过，骨架上的氧原子占据了
水分子的位置。

\textbf{20.} 骨架带负电，只交换阳离子；一个
带长链的大阴离子受到排斥，无法通过窗口，因此停留在
水中，继续发挥作用。

\textbf{21.} 脱水后，沸石强烈吸水进入其笼中；只有
小于窗口的分子才能进入，因此它按尺寸分离分子：
即分子尺度上的分子筛。

\textbf{22.} 较大的 $\ce{K+}$ 离子位于窗口中并使其变窄：只有较小的
分子（如水）才能通过。

\textbf{23.} 细长、线性的\textit{对位}异构体能嵌入孔道并快速扩散；
较宽的\textit{邻位}与\textit{间位}异构体扩散缓慢，停留并发生异构化。

\textbf{24.} 无水沸石 A 理论上每克可吸收 \qty{3.52}{mmol} 的 $\ce{Ca^2+}$。
\end{solution}
